% ============================================================
%  SUPPLEMENT 07 — STATISTICAL DEPENDENCE AND INFERENCE
%  Overlapping Windows, Block Bootstrap, Effective N, and Clustered SE
%  Companion to: ECDR_Data_Analysis.tex
% ============================================================
\documentclass[11pt, a4paper]{article}

\usepackage{fontspec}
\setmainfont{Noto Serif}
\setsansfont{Noto Sans}
\setmonofont{Noto Sans Mono}[Scale=0.88]

\usepackage[english]{babel}
\usepackage{amsmath, amssymb, amsthm}
\usepackage{mathtools}
\usepackage[margin=2.5cm]{geometry}
\usepackage{booktabs, tabularx, array, longtable}
\usepackage{xcolor}
\newcolumntype{Y}{>{\raggedright\arraybackslash}X}
\newcolumntype{P}[1]{>{\raggedright\arraybackslash}p{#1}}
\usepackage{hyperref}
\usepackage{enumitem}
\usepackage{microtype}
\usepackage{parskip}

\definecolor{ink}{RGB}{20,20,20}
\definecolor{teal}{RGB}{0,120,120}
\definecolor{muted}{RGB}{90,90,90}

\hypersetup{colorlinks=true, linkcolor=teal, urlcolor=teal, citecolor=teal}

\theoremstyle{definition}
\newtheorem{definition}{Definition}[section]
\newtheorem{example}{Example}[section]
\newtheorem{remark}{Remark}[section]
\theoremstyle{plain}
\newtheorem{theorem}{Theorem}[section]
\newtheorem{proposition}{Proposition}[section]

\newcommand{\R}{\mathbb{R}}
\newcommand{\E}{\mathbb{E}}
\newcommand{\Cov}{\operatorname{Cov}}
\newcommand{\Var}{\operatorname{Var}}
\newcommand{\corr}{\operatorname{corr}}
\newcommand{\code}[1]{\texttt{#1}}

\usepackage{titlesec}
\titleformat{\section}{\large\bfseries\color{teal}}{}{0em}{\thesection\quad}
\titleformat{\subsection}{\normalsize\bfseries\color{ink}}{}{0em}{\thesubsection\quad}

\begin{document}

\begin{center}
  {\LARGE\bfseries\color{teal} Statistical Dependence and Rigorous Inference}\\[0.5em]
  {\large\color{muted} Overlapping Windows, Block Bootstrap, Effective $N$, and Clustered Standard Errors}\\[1.2em]
  {\normalsize Supplement~07 to \textit{ECDR Data Analysis} \textbar{} October~2026}\\[0.3em]
  {\small\color{muted} This document is supplementary only. The main document is not modified.}
\end{center}

\vspace{1em}\hrule\vspace{1em}

\tableofcontents
\clearpage

% ============================================================
\section{The i.i.d. Assumption and Why It Fails Here}

Classical statistical inference (confidence intervals, hypothesis tests,
standard errors) assumes that observations $Z_1, Z_2, \ldots, Z_N$ are
\textbf{independent and identically distributed (i.i.d.)}: each observation is
drawn independently from the same distribution.

In the HAPT sliding-window dataset, this assumption fails at \emph{two levels}:

\subsection{Level 1: Window Overlap}

Adjacent windows share 50\% of their raw coordinates. Specifically, window $k$
contains raw samples at positions $[kS, kS + 128)$ and window $k+1$ at
$[(k+1)S, (k+1)S + 128)$. With stride $S = 64$, the overlap is:
\[
  \text{Overlap} = [kS + 64, kS + 128) = [(k+1)S, (k+1)S + 64).
\]
This is $128 - 64 = 64$ samples per channel, $64 \times 6 = 384$ values in total.
Out of 768 coordinates in each window, 384 are shared between consecutive windows.

\textbf{Consequence:} Any random variable $Z$ computed per-window (e.g., the $L_2$
distance from the mean, or the model logit output) has strong covariance between
adjacent windows, because they share half their input data.

\subsection{Level 2: Intra-Subject Clustering}

Multiple sessions from the same participant share that person's gait dynamics,
body mass, phone attachment angle, and movement style. A classifier that happens
to perform well on a particular subject's unique gait style will show inflated
accuracy within that subject's windows — not because of general robustness but
because of subject-specific memorization.

The 4 validation subjects (27--30) and 9 test subjects (2, 4, 9, 10, 12, 13,
18, 20, 24) provide subject-level independence by design. However, within each
subject, sessions share the subject's characteristics.

% ============================================================
\section{The True Variance of a Sample Mean}

\subsection{General Formula}

For a scalar summary statistic $Z_i$ computed per window $i$, the sample mean is:
\[
  \bar{Z} = \frac{1}{N} \sum_{i=1}^N Z_i.
\]
The true sampling variance (the variance of $\bar{Z}$ over hypothetical reruns
of the data collection experiment) is:
\begin{equation}\label{eq:true_var}
  \Var(\bar{Z})
  = \frac{1}{N^2} \sum_{i=1}^N \sum_{j=1}^N \Cov(Z_i, Z_j)
  = \frac{\sigma^2}{N} + \frac{1}{N^2} \sum_{i \ne j} \Cov(Z_i, Z_j),
\end{equation}
where $\sigma^2 = \Var(Z_i)$ (assuming equal marginal variances).

The naive i.i.d. standard error formula $\hat{\sigma}/\sqrt{N}$ assumes that
the double sum reduces to $N\sigma^2$ (all off-diagonal covariances are zero).
When covariances are \emph{positive} (adjacent windows are similar), the double sum
exceeds $N\sigma^2$, and the true standard error is \emph{larger} than $\hat{\sigma}/\sqrt{N}$.

\subsection{Quantifying the Inflation}

For a simple AR(1) (first-order autoregressive) model with lag-1 autocorrelation $\rho_1$:
\[
  \Cov(Z_i, Z_{i+k}) = \sigma^2 \rho_1^k.
\]
The variance of $\bar{Z}$ for a single session of length $T$ windows is:
\begin{align}
  \Var(\bar{Z}) &= \frac{\sigma^2}{T} \cdot \frac{1 + \rho_1}{1 - \rho_1}
  + O(T^{-2}) \quad \text{as } T \to \infty.
\end{align}
The inflation factor is $(1 + \rho_1)/(1 - \rho_1)$.

\begin{example}[AR(1) Inflation for Window Channel Mean]
  The channel-mean autocorrelation (Metric 4 in the data analysis) is approximately
  $\rho_1 \approx 0.97$ for \code{acc\_x}. The inflation factor:
  \[
    \frac{1 + 0.97}{1 - 0.97} = \frac{1.97}{0.03} \approx 65.7.
  \]
  The true standard error of the sample mean is $\sqrt{65.7} \approx 8.1$ times
  larger than the naive i.i.d. formula suggests. Reporting a 95\% confidence
  interval using the naive formula would be drastically overconfident.
\end{example}

% ============================================================
\section{Effective Sample Size: A Useful but Limited Concept}

\subsection{Definition}

The \textbf{effective sample size} $N_{\text{eff}}$ is defined as the number of
hypothetical i.i.d. observations that would give the same variance as the actual
$N$ correlated observations:
\begin{equation}
  N_{\text{eff}} = \frac{\sigma^2}{\Var(\bar{Z})} \cdot N
  \;=\; \frac{N}{1 + 2 \sum_{k=1}^{\infty} \rho_k},
\end{equation}
where $\rho_k = \Cov(Z_i, Z_{i+k})/\sigma^2$ is the autocorrelation at lag $k$.
For the AR(1) case: $N_{\text{eff}} = N(1 - \rho_1)/(1 + \rho_1)$.

\subsection{Why $N_{\text{eff}}$ Cannot Be Computed from Geometric Overlap}

A tempting shortcut: since windows overlap by 50\%, declare $N_{\text{eff}} = N/2$.
This is wrong for two reasons:

\begin{enumerate}[leftmargin=1.5em]
  \item \textbf{$N_{\text{eff}}$ depends on the statistic, not just the overlap.}
        For a statistic $Z_i$ that depends mostly on the first 64 samples (the
        non-overlapping part), adjacent windows are nearly independent and
        $N_{\text{eff}} \approx N$. For a statistic that depends on all 128 samples,
        the overlap inflates correlation. The 50\% geometric overlap does not
        mechanically translate to $N_{\text{eff}} = N/2$.
  \item \textbf{Non-linear classifier outputs break the formula.}
        The autocorrelation formula above is derived for linear statistics.
        Neural network predictions $\arg\max F_j(Z_i)$ are discrete, non-linear,
        and can exhibit clustering effects (the same label over many consecutive
        windows) that are not captured by second-order autocorrelation. The
        formula for $N_{\text{eff}}$ via $\sum \rho_k$ requires a stationary,
        weakly-dependent process — not guaranteed for a non-linear classifier on
        a non-stationary activity stream.
  \item \textbf{Subject clustering is separate from window overlap.}
        Even correcting for window overlap, windows from the same subject are
        correlated at a higher level. $N_{\text{eff}}$ for within-subject
        statistics is much smaller than $N_{\text{eff}}$ for cross-subject statistics.
\end{enumerate}

% ============================================================
\section{Block Bootstrap: The Right Tool for Dependent Data}

\subsection{Motivation}

The standard (i.i.d.) bootstrap resamples individual observations with replacement.
For time series data, this destroys the temporal dependence structure: resampled
sequences are not valid time series.

The \textbf{block bootstrap} resamples contiguous blocks of observations, preserving
local temporal dependence within each block.

\subsection{Subject-Level Block Bootstrap}

For the HAPT dataset, the appropriate block structure is \textbf{subject-level}:
each subject's collection of windows is treated as one block.

\textbf{Procedure for a single bootstrap replicate of test accuracy:}
\begin{enumerate}[leftmargin=1.5em]
  \item Let there be $S = 9$ test subjects.
  \item Sample $S$ subjects with replacement from the $S$ test subjects: some
        subjects appear multiple times, some not at all.
  \item Pool all windows from the sampled subjects.
  \item Compute the desired statistic (e.g., classification accuracy) on the pooled windows.
  \item Repeat steps 1--4 $B = 10{,}000$ times to obtain a bootstrap distribution.
  \item Report the 2.5th and 97.5th percentiles as a 95\% confidence interval.
\end{enumerate}

\subsection{Within-Subject Block Bootstrap}

If within-subject statistics are needed (e.g., session-level accuracy), a finer
block bootstrap can resample sessions (each session is a contiguous block) with
replacement within a subject. Within each session, windows are kept in order to
preserve temporal dependence.

\subsection{Practical Limitation: 4 Validation Subjects}

The validation split contains only $S_{\text{val}} = 4$ subjects (27, 28, 29, 30).
The subject-level block bootstrap has $4^4 = 256$ distinct bootstrap resamples
(compared to $9^9 = 387,420,489$ for the test set). Confidence intervals derived from
only 4 subjects are extremely wide and should not be interpreted as precise estimates.
This is a fundamental limitation of the split design, not a computational artifact.

% ============================================================
\section{Clustered Standard Errors}

\subsection{Definition}

Clustered standard errors are a regression-based alternative to block bootstrap.
For a statistic $Z_i$ from cluster (subject) $s(i)$, the \textbf{cluster-robust
variance estimator} is:
\begin{equation}
  \hat{V}_{\text{cluster}} = \frac{1}{N^2} \sum_{s=1}^{S} \left(\sum_{i: s(i) = s} (Z_i - \bar{Z})\right)^2.
\end{equation}
This is equivalent to summing the \emph{squared cluster totals} of residuals.
It allows arbitrary within-cluster dependence (not just AR(1)) and is consistent
as $S \to \infty$.

\subsection{Limitations}

\begin{itemize}[leftmargin=1.5em]
  \item Consistency requires $S \to \infty$. With 4 validation subjects, the
        asymptotic approximation is poor. Use bootstrap instead.
  \item Clustered SEs assume clusters are independent. Test subjects are
        independent across subjects (by split design), so this holds.
  \item For binary statistics (correct/incorrect per window), use clustered logistic
        regression or the bootstrap — not ordinary least squares clustered SEs.
\end{itemize}

% ============================================================
\section{What Evaluation Reports Must State}

Based on the statistical analysis, any evaluation of the ECDR system or the
reference classifier on HAPT data must explicitly state:

\begin{longtable}{@{}lp{9cm}@{}}
\toprule
Requirement & Explanation \\
\midrule
Population denominator & Number of subjects (not windows) contributing to the estimate. \\
Overlap handling & Whether results are reported on all windows or on independent
                    (non-overlapping) subsets. \\
Uncertainty quantification & Bootstrap CI or clustered SE; not naive $\pm 1.96\hat{\sigma}/\sqrt{N}$. \\
Activity-stratified results & Accuracy broken down by class, especially for postural transitions
                               with $< 100$ windows. \\
Split used & Train/val/test clearly identified; no mixing. \\
Evaluation mode & Pure windows only, all windows, or mixed stream — with explicit rationale. \\
\bottomrule
\caption{Required disclosures for statistically valid ECDR evaluation reports.}
\end{longtable}

% ============================================================
\section{The Normalization Population Bias}

The channel-wise standardization parameters $\mu_c, \sigma_c$ are fitted on the
5,693 pure 6-class training windows. These windows were constructed from overlapping
sliding windows of the training subjects' raw recordings.

\textbf{Bias mechanism:} Each raw sample $r[n]$ within a training recording appears
in multiple windows (up to $\lceil W/S \rceil = 2$ overlapping windows). Fitting
the mean $\mu_c$ as an average over all window samples implicitly weights each
raw sample by the number of windows it appears in. Samples near the middle of a
long continuous annotation (appearing in 2 windows) are weighted twice as heavily
as samples near annotation boundaries (appearing in 1 window).

\textbf{Magnitude of bias:} For the training set, the multiplicity distribution
(verified in the data analysis) shows that most samples appear in exactly 2 windows.
The induced bias in $\mu_c$ is:
\[
  \hat{\mu}_c^{\text{overlap}} - \mu_c^{\text{unique}} = \frac{\sum_n m_n r[n,c] / \sum_n m_n - \sum_n r[n,c] / T}{}
  \approx \frac{\text{boundary effect}}{T},
\]
where $m_n \in \{1, 2\}$ is the per-sample multiplicity. Because $\sum_n m_n / T \approx 2$
for most samples, the bias is small (a second-order effect at annotation boundaries).
However, if future work uses a different stride $S$, the multiplicity changes and
the bias should be re-evaluated.

\begin{thebibliography}{9}
\bibitem{main}
  ECDR Research Workspace. \textit{ECDR Data Analysis} (English edition).
  \code{ECDR\_Data\_Analysis.tex}, October~2026.
\bibitem{politis1994}
  D.N. Politis, J.P. Romano.
  \textit{The Stationary Bootstrap}.
  Journal of the American Statistical Association, 89(428):1303--1313, 1994.
\bibitem{cameron2011}
  A.C. Cameron, D.L. Miller.
  \textit{Robust Inference with Clustered Data}.
  In B.H. Baltagi (ed.), Econometrics of Panel Data. Springer, 2011.
\bibitem{efron1994}
  B. Efron, R.J. Tibshirani.
  \textit{An Introduction to the Bootstrap}. Chapman \& Hall, 1994.
\end{thebibliography}

\end{document}
